\newpage
\Solution{2}
We will begin with the second part of the question. The flipping is performed uniformly:
\begin{equation*}
	\boxed{\omega = \frac{\pi}{\tau}}
	\tag{$1$}
	\label{eq:j17_p2_1}
\end{equation*}

\begin{center}
	\begin{tikzpicture}
		% Ellipses
		\draw[line] (3, 4) ellipse (2cm and 0.25cm);
		\draw[line, rotate around={-30:(9, 4)}] (9, 4) ellipse (2cm and 0.25cm);
		\node[point] at (3, 4) {};
		\node[point] at (9, 4) {};
		
		% B-fields
		\draw[line, -{Stealth}] (3, 4) -- (3, 6);
		\draw[line, -{Stealth}] (9, 4) -- (9, 6);
		
		\node[above] at (3, 6) {$\vv{B}$};
		\node[above] at (9, 6) {$\vv{B}$};
		
		% Normals
		\draw[line, -{Stealth}] (3, 4) -- (3, 5);
		\draw[line, -{Stealth}, rotate around={-30:(9, 4)}] (9, 4) -- (9, 5);
		
		\node[right] at (3, 5) {$\hat{n}$};
		\node[right] at (9.5, 4.85) {$\hat{n}$};
		
		% Angle
		\coordinate (A) at (9, 6);
		\coordinate (B) at (9, 4);
		\coordinate (C) at (9.5, 4.85);
		
		\pic [draw, line, angle radius=0.5cm] {angle=C--B--A};
		
		\node[left] at (9.05, 4.61) {$\varphi$};
	\end{tikzpicture}
\end{center}

We can directly calculate the induced EMF in the ring as a function of time:
\begin{equation*}
	\varepsilon 
	= 
	-\frac{d\Phi_B}{dt}
	=
	- \frac{d}{dt}
	\left(
	\vv{B} \cdot \vv{S}
	\right)
	=
	-\frac{d}{dt}
	\left(
	\pi r^2 B \cos(\varphi)
	\right)
	=
	\pi r^2 B \sin(\varphi) \frac{d\varphi}{dt}
	\quad\Rightarrow\quad
	\boxed{
	\varepsilon(t)
	=
	\pi r^2 \omega B \sin(\omega t)
	}
	\tag{$2$}
	\label{eq:j17_p2_2}
\end{equation*}
Using equation (\ref{eq:j17_p2_1}) and Ohm's law:
\begin{equation*}
	i(t) = \frac{\varepsilon(t)}{R}
	\quad\Rightarrow\quad
	\boxed{i(t) = \frac{\pi^2 r^2 B}{R\tau} \sin\left(\frac{\pi t}{\tau}\right)}
	\tag{$3$}
	\label{eq:j17_p2_3}
\end{equation*}
The charge is then given by:
\begin{equation*}
	\Delta q
	=
	\int_0^\tau i(t)\,dt
	=
	\frac{1}{R}
	\int_0^\tau \varepsilon(t)\,dt
	=
	\frac{\pi r^2 B}{R}
	\int_0^\pi \sin\varphi\,d\varphi
	\quad\Rightarrow\quad
	\boxed{
		\Delta q 
		= 
		\frac{2\pi r^2 B}{R} 
		\approx 1{,}57\,\mathrm{\mu C}
		}
	\tag{$4$}
	\label{eq:j17_p2_4}
\end{equation*}
Note that this result does not depend on how the ring is flipped, provided it starts and ends in the same orientations. Consequently, the exact time dependence of the current  $i(t)$ is irrelevant for the total charge transferred. Only the initial and final magnetic fluxes matter:
\begin{equation*}
	\varepsilon = R \frac{dq}{dt} = - \frac{d\Phi_B}{dt}
	\quad\Rightarrow\quad
	\left|\Delta q \right|
	= 
	\frac{\left| \Delta \Phi_B \right|}{R} = 
	\frac{\pi r^2 B - (- \pi r^2 B)}{R}
	\quad\Rightarrow\quad
	\boxed{
		\left| \Delta q \right|
		= 
		\frac{2\pi r^2 B}{R} 
		\approx 1{,}57\,\mathrm{\mu C}
	}
\end{equation*}